\documentclass[11pt,a4paper]{article}
\usepackage[T1]{fontenc}
\usepackage[utf8]{inputenc}
\usepackage{lmodern}
\usepackage{amsmath,amssymb,amsthm,mathtools,mathrsfs}
\usepackage{graphicx,float,xcolor,multicol}
\usepackage[shortlabels]{enumitem}
\usepackage[margin=27mm]{geometry}
\usepackage{microtype,needspace,placeins}
\usepackage{tikz,tikz-cd}
\usetikzlibrary{arrows.meta,calc,positioning,patterns,decorations.pathreplacing}
\usepackage[colorlinks=true,linkcolor=teal,urlcolor=teal,citecolor=teal]{hyperref}
\setlength{\parindent}{0pt}
\setlength{\parskip}{.6em}
\setcounter{secnumdepth}{0}
\allowdisplaybreaks
\raggedbottom
\providecommand{\cross}{\times}
\DeclareUnicodeCharacter{2020}{\ensuremath{\dagger}}
\DeclareMathOperator{\supp}{supp}
\DeclareMathOperator*{\esssup}{ess\,sup}
\providecommand{\chapter}{\section}
\newenvironment{tool}[2]{\subsection*{Tool #1 --- #2}}{}
\newcommand{\ToolRef}[1]{Tool #1}
\newcommand{\FigureTag}[2]{\textit{Figure #1}}
\newcommand{\Result}[1]{\subsection*{#1}}
\newcommand{\Application}[1]{\subsubsection*{Application\if\relax\detokenize{#1}\relax\else\space--- #1\fi}}
\newcommand{\ProofHeading}{\par\textit{Proof.}\space}
\newcommand{\ProofOf}[1]{\par\textit{Proof of #1.}\space}
\newcommand{\RemarkHeading}{\par\textbf{Remark.}\space}
\newcommand{\NoteHeading}{\par\textbf{Note.}\space}
\newcommand{\ExerciseHeading}{\par\textbf{Exercise.}\space}

% Rendering support only; mathematical text is in each independent post.
\usepackage{booktabs,array,longtable}
\definecolor{HandoutBlue}{HTML}{2F6690}
\definecolor{HandoutNavy}{HTML}{17365D}
\newtheorem{theorem}{Theorem}
\newtheorem{lemma}[theorem]{Lemma}
\newtheorem{proposition}[theorem]{Proposition}
\newtheorem{corollary}[theorem]{Corollary}
\newtheorem{conjecture}[theorem]{Conjecture}
\newtheorem{definition}{Definition}
\newtheorem{example}{Example}
\newtheorem{namedtoolcounter}{Tool}
\newenvironment{namedtool}[1]{\begin{namedtoolcounter}[#1]}{\end{namedtoolcounter}}
\newtheorem*{claim}{Claim}
\newtheorem*{remark}{Remark}
\newtheorem*{exercise}{Exercise}
\newtheorem*{question}{Question}
\newtheorem*{observation}{Observation}
\newcommand{\SourceChapter}[2]{\section*{#2}}
\newcommand{\Lecture}[2]{\section*{Lecture #1: #2}}
\newcommand{\unclear}[1]{\textit{[unclear: #1]}}
\newcommand{\sourcegap}{\textit{[unfinished in the source]}}
\DeclareMathOperator{\dist}{dist}
\DeclareMathOperator{\id}{id}
\DeclareMathOperator{\Hol}{Hol}
\DeclareMathOperator{\Per}{Per}
\DeclareMathOperator{\Fix}{Fix}
\DeclareMathOperator{\diam}{diam}
\DeclareMathOperator{\Var}{Var}
\DeclareMathOperator{\Leb}{Leb}
\DeclareMathOperator{\Hom}{Hom}
\DeclareMathOperator{\End}{End}
\DeclareMathOperator{\Ann}{Ann}
\DeclareMathOperator{\Spec}{Spec}
\DeclareMathOperator{\Max}{Max}
\DeclareMathOperator{\Ob}{Ob}
\DeclareMathOperator{\Tor}{Tor}
\DeclareMathOperator{\Ass}{Ass}
\DeclareMathOperator{\Mor}{Mor}
\DeclareMathOperator{\Fun}{Fun}
\DeclareMathOperator{\Nat}{Nat}
\DeclareMathOperator{\Cone}{Cone}
\DeclareMathOperator{\MCG}{MCG}
\DeclareMathOperator{\Aut}{Aut}
\DeclareMathOperator{\Deck}{Deck}
\DeclareMathOperator{\SL}{SL}
\DeclareMathOperator{\PSL}{PSL}
\DeclareMathOperator{\Area}{Area}
\DeclareMathOperator{\im}{im}
\DeclareMathOperator{\PSh}{PSh}
\newcommand{\CRing}{\mathsf{CRings}}
\newcommand{\AMod}{A\text{-}\mathsf{Mod}}
\newcommand{\Ab}{\mathsf{Ab}}
\newcommand{\Z}{\mathbb Z}
\newcommand{\Q}{\mathbb Q}
\newcommand{\R}{\mathbb R}
\newcommand{\C}{\mathbb C}
\newcommand{\N}{\mathbb N}
\newcommand{\kk}{\mathbb k}
\renewcommand{\aa}{\mathfrak a}
\newcommand{\bb}{\mathfrak b}
\newcommand{\pp}{\mathfrak p}
\newcommand{\qq}{\mathfrak q}
\newcommand{\mm}{\mathfrak m}
\newcommand{\nn}{\mathfrak n}
\newcommand{\rad}{\mathrm r}
\newcommand{\cat}[1]{\mathsf{#1}}
\setcounter{secnumdepth}{0}
\allowdisplaybreaks

\graphicspath{{figures/miscellaneous/}}
\begin{document}
\section*{Small Gaps between Primes}
% BLOG-CONTENT-BEGIN
\begin{center}\small S. Viswanathan\end{center}
% Source page 1
Let $\mathbb P$ denote the set of primes. Enumerate them in increasing
order as $p_1,p_2,p_3,\ldots,p_n,\ldots$, where $p_n$ denotes the $n$th prime.
We ask a very basic question about the additive distribution of $\mathbb P$.
Fix a width parameter $h$. How many primes can you find in the interval
$[x,x+h]$, and how often?

It would suffice to study $H_m:=\liminf_{n\to\infty}(p_{n+m}-p_n)$.
Set $\widehat H_m:=\{p_{n+m}-p_n:p_n\in\mathbb P\}$.
Studying $\{H_m\}_m$, assuming that its terms are indeed numbers, is
extremely hard because of the additive nature of the problem.
Primes appear naturally only as factors, so they are suited for
multiplicative studies. Regardless, there are deep conjectures about this.
Of them all, the
% Source page 2
famous, perhaps even mythical, one is the Twin Prime Conjecture.
\setcounter{theorem}{17}\begin{conjecture}[Twin Prime Conjecture]
$H_1=2$, or equivalently, there are infinitely many twin primes.
\end{conjecture}
Recently, by deploying advanced sieve-theoretic tools, highly nontrivial
bounds for $H_m$ have been found, establishing that they are indeed numbers!
The first breakthrough came from Goldston, Pintz, and Y{\i}ld{\i}r{\i}m,
who developed a sieve technique to find at least two primes in an interval
$[a,b]$.
\paragraph{2005: Goldston, Pintz, and Y{\i}ld{\i}r{\i}m.}
They proved
\[
 \liminf_{n\to\infty}\frac{p_{n+1}-p_n}{\log p_n}=0.
\]
In fact, they proved the stronger version
$p_{n+1}-p_n\le(\log p_n)^{1/2+\varepsilon}$ infinitely often.
\paragraph{2013: Yitang Zhang.}
$\liminf_{n\to\infty}(p_{n+1}-p_n)=H_1\le70{,}000{,}000$.
\paragraph{2013: Polymath.}
$H_1\le4{,}680$. Tao organised a Polymath project, inviting mathematicians
from all over to optimise Zhang's technique.
It turns out that Zhang, and subsequently the Polymath group, used
Deligne's theorem, which proves the Riemann hypothesis for finite fields.
% Source page 3
A postdoc named James Maynard found a shortcut and surprised everyone.
He used elementary sieve methods.
\paragraph{2013: J. Maynard, independently.}
$H_1\le600$. Subsequently, Maynard and Tao's group joined forces.
\paragraph{2014: Maynard and Polymath.}
$H_1\le246$.
They were also able to prove the following bound for $H_m$.
\paragraph{2015: Maynard and Tao's group.}
$H_m\ll m^3e^{4m}$, asserting for the first time that $\{H_m\}_m$ is a
sequence of integers.
% Author check: the historical narrative and dated claims are reproduced, not updated.
This talk is about the 2005 work of Goldston, Pintz, and Y{\i}ld{\i}r{\i}m.
\setcounter{theorem}{18}\begin{conjecture}
$\displaystyle\liminf_{n\to\infty}\frac{p_{n+1}-p_n}{\log p_n}=0$.
\end{conjecture}
Erd\H{o}s was the first to make nontrivial progress. He proved
$\liminf_{n\to\infty}(p_{n+1}-p_n)/\log p_n\le1$.
This was improved to $0.24$. Finally, in 2005, Goldston, Pintz, and
Y{\i}ld{\i}r{\i}m proved the following.
% Source page 6
\setcounter{theorem}{19}\begin{theorem}[Goldston, Pintz, and Y{\i}ld{\i}r{\i}m]
\label{thm:m5-gpy}
$\displaystyle\liminf_{n\to\infty}\frac{p_{n+1}-p_n}{\log p_n}=0$.
\end{theorem}
The quantity $\log p$ is usually called the \emph{average gap} because,
according to the Prime Number Theorem, there is at least one prime on
average in each interval
$(k\log n,(k+1)\log n)$, where $0\le k\le n/\log n$.
With Theorem~\href{https://codezen1729.github.io/math-blog/\#/post/small-gaps-between-primes}{20}, one can ask whether
$\liminf_{p\to\infty}(p_{\mathrm{next}}-p)<\infty$, that is, whether
bounded gaps between primes exist.

\setcounter{theorem}{20}\begin{conjecture}
$\displaystyle\limsup_{p\to\infty}
\frac{p_{\mathrm{next}}-p}{(\log p)^2}=1$.
\end{conjecture}
Consider the sequence of consecutive numbers $m!+2,m!+3,\ldots,m!+m$.
If $m!\sim x$, then the length of $(m!+2,m!+m)$ is
$\sim\log x/\log\log x$, which is far less than $(\log x)^2$.
In fact, it is dominated even by the average gap.

% Source page 7
Regardless, one must be ambitious. Goldston, Pintz, and Y{\i}ld{\i}r{\i}m
were! They wanted to prove $\boxed{H_1\le16}$.
Their idea is rather benign.
Let $\mathcal H=\{h_1,h_2,\ldots,h_k\}$, where $h_1<h_2<\cdots<h_k$.
Suppose one can find a sequence $\{a(n)\}_n$ of nonnegative numbers such
that, for every $j=1,2,\ldots,k$,
\begin{equation}
 \sum_{\substack{x\le n\le2x\\n+h_j\in\mathbb P}}a(n)
 >\frac1k\sum_{x\le n\le2x}a(n).
 \tag{$\star$}\label{eq:m5-star}
\end{equation}
Then
\[
 \sum_{j=1}^k\left(\sum_{\substack{x\le n\le2x\\n+h_j\in\mathbb P}}a(n)\right)
 =\sum_{x\le n\le2x}\#\{1\le j\le k:n+h_j\in\mathbb P\}\,a(n)
 >\sum_{x\le n\le2x}a(n).
\]
Hence $\#\{1\le j\le k:n+h_j\in\mathbb P\}\ge2$, giving
$p_{n+1}-p_n\le\diam(\{h_j\}_{j=1}^k)$.
Inspired by Selberg, their choice of $\{a(n)\}_n$ is as follows.
The idea is to choose $a(n)=1$ if the numbers
$\{n+h_j\}_{j=1}^k$ are all prime, because, even without
\href{https://codezen1729.github.io/math-blog/\#/post/small-gaps-between-primes}{($\star$)}, we would at least have an upper bound on the number
of prime $k$-tuples.
% Source page 8
Let $\{\lambda_d\}_d$ be a sequence of real numbers such that
$\lambda_1=1$ and $\lambda_d=0$ for all $d>R$. Set
\[
 a(n):=\left(\sum_{d\mid(n+h_1)\cdots(n+h_k)}\lambda_d\right)^2.
\]
If $R<x\le n$ and all the $n+h_j$ are prime, then $a(n)=1$.
Expanding the square, one has
\[
 \sum_{d_1,d_2}\lambda_{d_1}\lambda_{d_2}
 \left(\sum_{\substack{x\le n\le2x\\
 d_1\mid(n+h_1)\cdots(n+h_k)\\d_2\mid(n+h_1)\cdots(n+h_k)}}1\right)
 =
 \sum_{d_1,d_2}\lambda_{d_1}\lambda_{d_2}
 \left(\sum_{\substack{x\le n\le2x\\{}[d_1,d_2]\mid(n+h_1)\cdots(n+h_k)}}1\right),
\]
since $[d_1,d_2]\mid m$ if and only if $d_1\mid m$ and $d_2\mid m$.
\setcounter{theorem}{21}\begin{theorem}[Assuming the Elliott--Halberstam conjecture]
$H_1\le16$.
\end{theorem}
% Source page 11
This failed attempt to get $H_1\le16$ led to decades of fascinating
progress, leading up to the following state-of-the-art theorem.
\setcounter{theorem}{22}\begin{theorem}[Polymath]
$H_1\le246$.
\end{theorem}
It was argued that one cannot beat $6$ using sieve-theoretic tools because
of something called the \emph{parity obstruction}. Thus this theory is,
at the moment, stagnant. Tao remarks that progress is one or two clever
ideas away!
% The historical "at the moment" and conditional hypothesis are retained as source content, not current claims.
% BLOG-CONTENT-END
\end{document}
